\documentclass[12pt]{article}
\usepackage[utf8]{inputenc}
\usepackage[T1]{fontenc}
\usepackage[margin=1in]{geometry}
\usepackage{setspace}
\usepackage{xcolor}
\usepackage[colorlinks=true, urlcolor=blue!60!black, linkcolor=blue!60!black]{hyperref}
\setlength{\parindent}{0em}
\setlength{\parskip}{0.45em}
\renewcommand{\baselinestretch}{1.0}
\begin{document}
\begin{center}
{\Large \textit{Griffin v.\ Breckenridge}, 403 U.S.\ 88 (1971)\par}
\vspace{0.3em}
{\normalsize Case Brief --- One-Page Edition\par}
\end{center}

\textbf{Memory Jogger.} Private racial conspiracies are actionable under \S~1985(3) without state action (\href{https://supreme.justia.com/cases/federal/us/403/88/}{403 U.S. 88}).

\textbf{Facts.} Two white Mississippians, believing that driver R. G. Grady was a civil rights worker, blocked his car on a highway, forced him and his Black passengers out at gunpoint, and clubbed them. Neither acted under color of law. The victims sued under the Ku Klux Klan Act's \S~1985(3).

\textbf{Procedural History.} The district court dismissed under \textit{Collins v.\ Hardyman}; the Fifth Circuit affirmed.

\textbf{Votes.} Reversed, 9--0 (8--1 on interstate travel). Stewart wrote for the Court; Harlan concurred in part.

\textbf{Holding.} Section 1985(3) reaches private conspiracies with racial or class-based invidious animus. Congress could reach this one under the Thirteenth Amendment and its power to protect interstate travel.

\textbf{Reasoning.} The statute's text (conspiring or going in disguise on the highway), structure, and history show it targets private Klan violence; the animus requirement keeps it from becoming a general federal tort law. Under \textit{Jones v.\ Alfred H.\ Mayer Co.}, Congress may remedy the badges and incidents of slavery, and interstate travel is protected against private interference. Harlan relied on the Thirteenth Amendment alone.

\textbf{Examples.} \emph{Same-Side:} a racist plot to drive out a Black family. \emph{Opposite-Side:} anti-union violence or clinic blockades. \emph{Fence-Sitter:} an attack on a religious minority.

\textbf{Critique.} Originalists say it stretches the Thirteenth Amendment; progressives blame its vague animus test for later narrowing.

\textbf{Key Quotation.} ``The language requiring intent to deprive of equal protection, or equal privileges and immunities, means that there must be some racial, or perhaps otherwise class-based, invidiously discriminatory animus behind the conspirators' action. The conspiracy, in other words, must aim at a deprivation of the equal enjoyment of rights secured by the law to all.'' --- Justice Stewart, Opinion of the Court, 403 U.S. at 102.

\end{document}
